\begin{ejer}
Demuestra que
\[
\left(k[x,y]/(xy)\right)_{\ubar x}\cong k[x]_x,
\qquad\text{con }y\longmapsto0
\]
\end{ejer}
\begin{sol}
La flecha $k[x,y]\to k[x]$ tal que $y\mapsto0$, inducida por la propiedad universal, se compone con el morfismo canónico $k[x]\to k[x]_x$, lo cual resulta en
\[
k[x,y]\longrightarrow k[x]_x,
\qquad
\p(x,y)\longmapsto\frac{\p(x,0)}{1},
\qquad
y\longmapsto\frac{0}{1}
\]
En particular, $xy\mapsto0$, así que su kernel contiene al ideal $(xy)$ e induce un único morfismo
\[
\varphi\colon\frac{k[x,y]}{(xy)}\longrightarrow k[x]_x,
\qquad
\ubar{\p(x,y)}\longmapsto\frac{\p(x,0)}{1},
\qquad
\ubar y\longmapsto\frac{0}{1}
\]
Como $\varphi(\ubar x)=\frac{x}{1}$ es invertible en $k[x]_x$, $\varphi$ induce a su vez
\[
\Phi\colon\left(\frac{k[x,y]}{(xy)}\right)_{\ubar x}\longrightarrow k[x]_x,
\qquad
\Phi\left(\frac{\ubar{\p(x,y)}}{\ubar x^n}\right)
=\frac{\varphi(\ubar{\p(x,y)})}{x^n}
\]
Construyamos su inversa: el resultado de componer morfismos canónicos
\[
k[x]\longrightarrow k[x,y]\longrightarrow\frac{k[x,y]}{(xy)}
\longrightarrow\left(\frac{k[x,y]}{(xy)}\right)_{\ubar x},
\qquad
\p(x)\longmapsto\frac{\ubar{\p(x)}}{1}
\]
manda $x$ al elemento invertible $\frac{\ubar x}{1}$, por lo que induce un morfismo
\[
\Psi\colon k[x]_x\longrightarrow\left(\frac{k[x,y]}{(xy)}\right)_{\ubar x},
\qquad
\Psi\left(\frac{\p(x)}{x^n}\right)=\frac{\ubar{\p(x)}}{\ubar x^n}
\]
Solo resta probar que $\Phi$ y $\Psi$ son inversos entre sí:
\[
(\Phi\circ\Psi)\left(\frac{\p(x)}{x^n}\right)
=\Phi\left(\frac{\ubar{\p(x)}}{\ubar x^n}\right)
=\frac{\varphi(\ubar{\p(x)})}{x^n}
=\frac{\p(x)}{x^n}
\]
Para la otra composición basta probar que
\[
(\Psi\circ\Phi)\left(\frac{\ubar{\p(x,y)}}{\ubar x^n}\right)
=\frac{\ubar{\p(x,y)}}{\ubar x^n}
\]
gracias a que $\ubar y=0$ en $\left(k[x,y]/(xy)\right)_{\ubar x}$, pues $\ubar x\ubar y=0$ implica $\ubar y=\frac{1}{\ubar x}\ubar x\ubar y=0$. Entonces
\[
(\Psi\circ\Phi)\left(\frac{\ubar{\p(x,y)}}{\ubar x^n}\right)
=\Psi\left(\frac{\varphi(\ubar{\p(x,y)})}{x^n}\right)
=\Psi\left(\frac{\p(x)}{x^n}\right)
=\frac{\ubar{\p(x)}}{\ubar x^n}
\]
\end{sol}